% Owner: Abhinav (P8.5)
\section{Introduction}\label{sec:introduction}

Indian cities are growing fast into the land around them. Bengaluru is a clear case: over four
decades its built-up area grew by 584\,\%, while vegetation fell by 66\,\% and water bodies by
74\,\% \citep{ramachandra2012}. Planners therefore keep facing the same question: \emph{where
can new urban development reasonably go?}

The classic GIS answer is a land suitability analysis. Factor maps such as slope, access to
roads and current land cover are rated, weighted and overlaid into one suitability surface
\citep{collins2001}. Since the 1990s the weights usually come from multi-criteria decision
analysis (MCDA), most often the Analytic Hierarchy Process \citep[AHP;][]{saaty1980}, in which
experts compare criteria in pairs \citep{malczewski2004}. The method is transparent and easy to
explain, but it has two well-known weaknesses. The weights are subjective, and changing only the
weighting method moves the boundaries of the suitable zones \citep{mosadeghi2015}. And the
resulting maps are rarely checked against what actually happened on the ground
\citep{malczewski2004, parry2018}.

Freely available satellite land use / land cover (LULC) data now make a different approach
possible. Annual 10\,m LULC maps \citep{karra2021} show where land has actually been built on,
year by year. A model can learn which surroundings preceded new development, and its predictions
can be tested against growth that happened later.

This project builds such an analysis for \textbf{Bengaluru South}, entirely in Python. It answers
three questions, set at the project meeting:
\begin{enumerate}
  \item \textbf{Land classification.} Which land is built-up, which is forest, and which is
        usable (potential) land? The land is grouped by clustering LULC fractions and terrain
        on a 100\,m grid.
  \item \textbf{Building similarity.} Given an existing building at location~A, which
        locations~B on usable land have similar surroundings (terrain, access to roads and
        water, nearby land cover, built-up density)?
  \item \textbf{Suitability.} How do data-driven scores (similarity to existing buildings and a
        random forest that learns from recent change) compare with a classic MCDA/AHP overlay
        when all of them are tested against real growth?
\end{enumerate}

The test is strict: every score is computed with data as of 2018 only and compared with growth
that persisted through 2022 and 2023. Rules against leakage of future information are enforced
in code. The result is a three-class land map, a graded suitability map with candidate sites, and
an honest comparison of the methods, including where they fail.

Section~\ref{sec:literature} reviews related work. Sections~\ref{sec:study-area}
and~\ref{sec:data} describe the study area and the data. The methods and the validation design
follow in Sections~\ref{sec:methodology} and~\ref{sec:validation_design}, the results in
Section~\ref{sec:results}, and a discussion and conclusion close the report.
